% Options for packages loaded elsewhere
\PassOptionsToPackage{unicode}{hyperref}
\PassOptionsToPackage{hyphens}{url}
\PassOptionsToPackage{dvipsnames,svgnames*,x11names*}{xcolor}
%
\documentclass[
  10pt,
]{article}
\usepackage{amsmath,amssymb}
\usepackage{lmodern}
\usepackage{iftex}
\ifPDFTeX
  \usepackage[T1]{fontenc}
  \usepackage[utf8]{inputenc}
  \usepackage{textcomp} % provide euro and other symbols
\else % if luatex or xetex
  \usepackage{unicode-math}
  \defaultfontfeatures{Scale=MatchLowercase}
  \defaultfontfeatures[\rmfamily]{Ligatures=TeX,Scale=1}
\fi
% Use upquote if available, for straight quotes in verbatim environments
\IfFileExists{upquote.sty}{\usepackage{upquote}}{}
\IfFileExists{microtype.sty}{% use microtype if available
  \usepackage[]{microtype}
  \UseMicrotypeSet[protrusion]{basicmath} % disable protrusion for tt fonts
}{}
\makeatletter
\@ifundefined{KOMAClassName}{% if non-KOMA class
  \IfFileExists{parskip.sty}{%
    \usepackage{parskip}
  }{% else
    \setlength{\parindent}{0pt}
    \setlength{\parskip}{6pt plus 2pt minus 1pt}}
}{% if KOMA class
  \KOMAoptions{parskip=half}}
\makeatother
\usepackage{xcolor}
\IfFileExists{xurl.sty}{\usepackage{xurl}}{} % add URL line breaks if available
\IfFileExists{bookmark.sty}{\usepackage{bookmark}}{\usepackage{hyperref}}
\hypersetup{
  pdftitle={CUAD Legal Document Analysis --- Complete Project Report},
  pdfauthor={Afnan Mazumdar (24141229) · Shoyeb Hasan Sayem (22101386) · Jerin Aktar (22101279); Supervisor: Utsha Kumar Roy},
  colorlinks=true,
  linkcolor={Maroon},
  filecolor={Maroon},
  citecolor={Blue},
  urlcolor={Blue},
  pdfcreator={LaTeX via pandoc}}
\urlstyle{same} % disable monospaced font for URLs
\usepackage[a4paper, margin=2cm]{geometry}
\usepackage{longtable,booktabs,array}
\usepackage{calc} % for calculating minipage widths
% Correct order of tables after \paragraph or \subparagraph
\usepackage{etoolbox}
\makeatletter
\patchcmd\longtable{\par}{\if@noskipsec\mbox{}\fi\par}{}{}
\makeatother
% Allow footnotes in longtable head/foot
\IfFileExists{footnotehyper.sty}{\usepackage{footnotehyper}}{\usepackage{footnote}}
\makesavenoteenv{longtable}
\setlength{\emergencystretch}{3em} % prevent overfull lines
\providecommand{\tightlist}{%
  \setlength{\itemsep}{0pt}\setlength{\parskip}{0pt}}
\setcounter{secnumdepth}{-\maxdimen} % remove section numbering
\usepackage{booktabs}
\usepackage{etoolbox}
\AtBeginEnvironment{longtable}{\small}
\renewcommand{\arraystretch}{1.15}
\setlength{\parskip}{4pt}
\usepackage{newunicodechar}
\newunicodechar{→}{\ensuremath{\rightarrow}}
\usepackage{fvextra}
\DefineVerbatimEnvironment{Highlighting}{Verbatim}{breaklines,commandchars=\\\{\}}
\fvset{breaklines=true,fontsize=\small}
\usepackage[htt]{hyphenat}
\ifLuaTeX
  \usepackage{selnolig}  % disable illegal ligatures
\fi

\title{CUAD Legal Document Analysis --- Complete Project Report}
\usepackage{etoolbox}
\makeatletter
\providecommand{\subtitle}[1]{% add subtitle to \maketitle
  \apptocmd{\@title}{\par {\large #1 \par}}{}{}
}
\makeatother
\subtitle{AI-Powered Legal Document Analysis System · BRAC University, Dept. of
CSE}
\author{Afnan Mazumdar (24141229) · Shoyeb Hasan Sayem (22101386) · Jerin Aktar
(22101279) \and Supervisor: Utsha Kumar Roy}
\date{Updated 2026-08-14}

\begin{document}
\maketitle

{
\hypersetup{linkcolor=}
\setcounter{tocdepth}{3}
\tableofcontents
}
A step-by-step account of everything done: data preparation, EDA, training of
all three models, the 80/20 train/test split, testing, and the evaluation
metrics (including the confusion matrix). Numbers below are the actual results
from our runs.

\begin{center}\rule{0.5\linewidth}{0.5pt}\end{center}

\subsection{0. Overview}

We built a three-model pipeline on the full CUAD dataset (510 contracts, 41
clause categories):

\begin{longtable}[]{@{}
  >{\raggedright\arraybackslash}p{(\columnwidth - 6\tabcolsep) * \real{0.25}}
  >{\raggedright\arraybackslash}p{(\columnwidth - 6\tabcolsep) * \real{0.25}}
  >{\raggedright\arraybackslash}p{(\columnwidth - 6\tabcolsep) * \real{0.25}}
  >{\raggedright\arraybackslash}p{(\columnwidth - 6\tabcolsep) * \real{0.25}}@{}}
\toprule
Stage & Model & Job & Test result (20\% held-out) \\
\midrule
\endhead
\textbf{2A} & Presence classifier & Which of the 41 clause types are in a
contract? & \textbf{micro-F1 0.776, Accuracy 85.49\%} \\
\textbf{2B} & Span extractor & Locate the exact clause text & \textbf{token-F1
0.763} \\
\textbf{1} & Summarizer & Rewrite a clause in plain English & \textbf{ROUGE-L
0.775} \\
\bottomrule
\end{longtable}

\begin{quote}
Numbers below are reproduced by the standalone \textbf{\texttt{test.ipynb}},
which loads the saved 80/20-trained models from disk and evaluates them on the
held-out 20\% test set.
\end{quote}

Everything below is the sequence of steps that produced these numbers.

\begin{center}\rule{0.5\linewidth}{0.5pt}\end{center}

\subsection{Step 1 --- Data source}

The official Atticus Project \texttt{theatticusproject/cuad} dataset on Hugging
Face. Its default \texttt{load\_dataset()} exposes only raw PDFs (not
trainable), so the annotated data was pulled from two repo files:

\begin{itemize}
\tightlist
\item
  \texttt{CUAD\_v1/master\_clauses.csv} --- 83-column wide table (used for EDA
  and label wrangling).
\item
  \texttt{CUAD\_v1/CUAD\_v1.json} --- SQuAD format with character offsets (used
  for full context + span training).
\end{itemize}

\subsection{Step 2 --- Data preparation (wide → long melt)}

\texttt{master\_clauses.csv} is wide: 41 categories x (clause-text column +
answer column). We:

\begin{itemize}
\tightlist
\item
  \textbf{Paired columns by position, not by name} --- the header
  \texttt{Notice\ Period\ To\ Terminate\ Renewal-\ Answer} has a stray space
  that breaks string matching; position pairing is robust to it.
\item
  \textbf{Defined ``present'' = non-empty clause list} (rather than the Yes/No
  answer), which works uniformly across the 8 entity categories (dates, parties)
  and hands us the clause spans for free.
\end{itemize}

\textbf{Result:} 510 contracts → \textbf{20,910 labeled (contract, category)
rows}, 32.1\% positive overall.

\subsection{Step 3 --- Exploratory Data Analysis (EDA)}

\begin{itemize}
\tightlist
\item
  \textbf{Severe class imbalance.} Positive rate ranges from
  \texttt{Document\ Name} 100\% and \texttt{Parties} 99.8\% down to
  \texttt{Source\ Code\ Escrow} 2.5\%. Median category \textasciitilde{} 23\%
  positive.
\item
  \textbf{The rare tail is near few-shot.} Some categories have only a handful
  of positives in the entire dataset.
\item
  \textbf{Contracts are long.} Median \textasciitilde{} 33,000 chars; max
  \textgreater{} 338,000 chars --- about 16x a transformer's input window.
\item
  \textbf{Clauses are short.} Median extracted clause \textasciitilde{} 196
  chars. So a clause fits easily in a window; the challenge is \emph{finding} it
  in a huge document.
\end{itemize}

\begin{center}\rule{0.5\linewidth}{0.5pt}\end{center}

\subsection{Step 4 --- The 80/20 train/test split}

Following the supervisor's instruction, the dataset is split \textbf{80/20 at
the contract level}:

\begin{longtable}[]{@{}llll@{}}
\toprule
& Contracts & (contract x category) rows & Positive rate \\
\midrule
\endhead
\textbf{Train (80\%)} & 408 & 16,728 & 32.0\% \\
\textbf{Test (20\%)} & 102 & 4,182 & 32.2\% \\
\textbf{Total} & 510 & 20,910 & 32.1\% \\
\bottomrule
\end{longtable}

\begin{itemize}
\tightlist
\item
  \textbf{Contract-level} (whole contracts go to one side) so no clause from a
  training contract leaks into test.
\item
  \textbf{Seeded} (\texttt{SEED=42}) and \textbf{reused across all three models}
  --- ``test'' means the same 102 unseen contracts everywhere.
\item
  Split files exported: \texttt{data/train\_80.csv}, \texttt{data/test\_20.csv}.
\end{itemize}

\textbf{Is the split actually balanced?} The aggregate positive rates (32.0\% vs
32.2\%) agree, but that alone could hide a per-category skew. Comparing all 41
categories directly (\texttt{analysis/ch3\_chart.py}):

\begin{longtable}[]{@{}
  >{\raggedright\arraybackslash}p{(\columnwidth - 2\tabcolsep) * \real{0.50}}
  >{\raggedright\arraybackslash}p{(\columnwidth - 2\tabcolsep) * \real{0.50}}@{}}
\toprule
Check & Value \\
\midrule
\endhead
Correlation between train and test positive rates & \textbf{r = 0.986} \\
Mean absolute difference across the 41 categories & \textbf{3.4 percentage
points} \\
Largest single gap & \textbf{14.0 pp} --- \emph{Termination For Convenience}
(47.1\% test vs 33.1\% train) \\
\bottomrule
\end{longtable}

The split is well balanced overall. We report the one real outlier rather than
smoothing it over: because \emph{Termination For Convenience} is
over-represented in test, it carries more weight in the micro-averaged score
than its training frequency implies. The split was drawn once at random and
never re-drawn.

\begin{center}\rule{0.5\linewidth}{0.5pt}\end{center}

\subsection{Step 5 --- Training Stage 2A: Presence classifier}

\subsubsection{5.1 Strong baseline first --- TF-IDF + Logistic Regression
(trained on 80\%)}

Classical ML has no length limit, so TF-IDF reads the whole contract. One
class-balanced Logistic Regression per category:

\begin{longtable}[]{@{}ll@{}}
\toprule
Metric & F1 \\
\midrule
\endhead
micro (every contract x category equal) & \textbf{0.775} \\
weighted (by support) & 0.777 \\
macro (all 41 categories) & 0.572 \\
macro (\textgreater=10 test positives, 34 cats) & 0.666 \\
\bottomrule
\end{longtable}

\subsubsection{5.2 The retrieval-ceiling check (why we did not window naively)}

Before training a transformer, we measured the \textbf{retrieval hit-rate} ---
can a query find the window that contains the gold clause? Best strategy
(category name, top-3) capped at \textbf{64\%}, i.e.~a retrieval-windowed
transformer's recall would be \emph{below} the baseline. Measuring this first
saved a doomed training run.

\subsubsection{5.3 The fix --- sliding-window max-pool (Multiple-Instance
Learning)}

We trained a \textbf{window-level DistilBERT}:
\texttt{(question,\ window)\ →\ does\ this\ window\ contain\ the\ clause?} At
inference we run \textbf{all} windows and \textbf{max-pool} (present if any
window fires), which removes the ceiling.

\begin{itemize}
\tightlist
\item
  Windowing: 2,000-char windows, 1,500-char stride → a 500-char overlap, which
  is wider than the median 196-char clause, so no clause can fall between two
  windows. A median-length contract yields \textbf{22 windows}; the corpus mean
  is \textbf{35} and the longest contract gives \textbf{226}.
\item
  Training examples: \textbf{53,662 windows} (39.2\% positive), subsampled to
  \textbf{24,000} for a laptop run.
\item
  Config: DistilBERT-base, 2 epochs, batch 16, LR 2e-5, max length 256.
\item
  Optimizer (read back from the saved \texttt{training\_args.bin}): AdamW,
  \textbf{linear decay, no warmup}, gradient-norm clipping at \textbf{1.0},
  3,000 total steps. Schedule plotted by \texttt{analysis/lr\_schedule.py}.
\item
  Training loss 0.319 → 0.261; validation loss 0.278 → 0.257;
  \textbf{\textasciitilde52 min} on Apple MPS.
\item
  Standalone test micro-F1: \textbf{0.694} (loses to the baseline ---
  max-pooling inflates false positives: the max is taken over every window, so
  one spurious window anywhere flips the whole contract).
\end{itemize}

\subsubsection{5.4 The winning model --- complementary AND-ensemble}

The two models are complementary (transformer wins on semantically subtle
categories, baseline on keyword-heavy ones). Combined with parameter-free rules:

\begin{longtable}[]{@{}ll@{}}
\toprule
Model & micro-F1 (20\% test) \\
\midrule
\endhead
\textbf{Ensemble AND (both must agree)} & \textbf{0.776} (best) \\
TF-IDF baseline & 0.775 \\
Ensemble AVG & 0.718 \\
Transformer (max-pool) alone & 0.694 \\
Ensemble OR & 0.696 \\
\bottomrule
\end{longtable}

The AND-rule filters each model's uncorrelated false positives → precision up →
best F1.

\subsubsection{5.5 How much of that ranking is statistically real?}

A 0.776-vs-0.775 margin on 102 contracts is one decision in a thousand, so we
tested it rather than asserting it. Method: \textbf{cluster bootstrap} ---
resample the 102 \emph{contracts} with replacement (not the 4,182 cells, which
are correlated within a contract), recompute micro-F1 for both models on the
same resample, 10,000 times. Script: \texttt{analysis/bootstrap\_ci.py},
\texttt{analysis/extra\_ci.py}.

\begin{longtable}[]{@{}
  >{\raggedright\arraybackslash}p{(\columnwidth - 8\tabcolsep) * \real{0.20}}
  >{\raggedright\arraybackslash}p{(\columnwidth - 8\tabcolsep) * \real{0.20}}
  >{\raggedright\arraybackslash}p{(\columnwidth - 8\tabcolsep) * \real{0.20}}
  >{\raggedright\arraybackslash}p{(\columnwidth - 8\tabcolsep) * \real{0.20}}
  >{\raggedright\arraybackslash}p{(\columnwidth - 8\tabcolsep) * \real{0.20}}@{}}
\toprule
Comparison & Difference in micro-F1 & 95\% CI & P(\textgreater0) & Verdict \\
\midrule
\endhead
AND-ensemble vs TF-IDF baseline & +0.0014 & {[}-0.0068, +0.0089{]} & 0.626 &
\textbf{tie --- CI contains zero} \\
TF-IDF baseline vs transformer & +0.0805 & {[}+0.0626, +0.0992{]} & 1.000 &
\textbf{real} \\
AND-ensemble vs transformer & +0.0818 & {[}+0.0654, +0.0987{]} & 1.000 &
\textbf{real} \\
\bottomrule
\end{longtable}

\textbf{The two findings are not on equal footing, and this is the honest
headline:}

\begin{itemize}
\tightlist
\item
  \textbf{The ensemble's win over the baseline is not statistically real.} The
  interval straddles zero and the ensemble leads in only 62.6\% of resamples.
  Individually: AND 0.776 (95\% CI {[}0.759, 0.792{]}), TF-IDF 0.775 (95\% CI
  {[}0.758, 0.791{]}) --- almost entirely overlapping. Read the top two rows of
  the table above as a \textbf{tie}.
\item
  \textbf{The transformer's loss to the baseline is real.} The baseline wins in
  \textbf{all 10,000} resamples.
\end{itemize}

This does not remove the reason for using the ensemble --- that reason was never
the 0.001 of micro-F1. The AND rule earns its place through the \textbf{error
profile} it produces (balanced precision and recall, see §8.1), not through its
rank.

\subsubsection{5.6 Two controls the bootstrap cannot run}

A bootstrap asks ``could this gap be sampling noise?'' It cannot ask ``was the
transformer handicapped?'' Both alternative explanations were tested
(\texttt{analysis/presence\_run.py}, \texttt{analysis/compare\_runs.py}).

\textbf{Control A --- training budget.} The baseline is fitted on all 408
contracts in full; the transformer saw 24,000 of 53,662 windows. Retrained on
\textbf{all 53,662} (140.7 min), everything else identical:

\begin{longtable}[]{@{}llll@{}}
\toprule
Transformer training windows & micro-F1 & Precision & Recall \\
\midrule
\endhead
24,000 (44.7\%) & 0.6943 & 0.5620 & 0.9080 \\
\textbf{53,662 (100\%)} & \textbf{0.6939} & 0.5546 & 0.9266 \\
\bottomrule
\end{longtable}

Difference \textbf{-0.0004}, CI {[}-0.0098, +0.0095{]}; gap to TF-IDF holds at
\textbf{+0.0809}. More data made the model fire \emph{more}, not better ---
recall up, precision down --- which is the max-pooling failure mode over
\textasciitilde35 windows. \textbf{The deficit is architectural, not a budget
artifact.}

\textbf{Control B --- random seed.} Three runs at the 24,000 budget:

\begin{longtable}[]{@{}lllllll@{}}
\toprule
Model & Seed 42 & Seed 43 & Seed 44 & Mean & SD & Range \\
\midrule
\endhead
Transformer & 0.6943 & 0.6928 & 0.7082 & 0.6984 & 0.0085 & 0.0154 \\
AND ensemble & 0.7761 & 0.7774 & 0.7803 & 0.7780 & 0.0022 & 0.0042 \\
\bottomrule
\end{longtable}

The ensemble's SD (0.0022) and range (0.0042) both \textbf{exceed its +0.0014
margin} --- reseeding moves the score further than the result being claimed for
it. The baseline's lead over the transformer is 9.5 SDs. The ensemble is also 4x
less seed-variable than the transformer alone, which is a better argument for it
than its rank.

\subsubsection{5.7 Does the AND rule punish rare categories?}

Splitting the 41 categories by training frequency
(\texttt{analysis/rare\_tail.py}):

\begin{longtable}[]{@{}lllllll@{}}
\toprule
Group & Model & Precision & Recall & F1 & TP & FN \\
\midrule
\endhead
\textbf{Rarest 10} (70 test pos.) & Transformer & 0.261 & 0.500 & 0.343 & 35 &
35 \\
& \textbf{AND ensemble} & 0.586 & 0.243 & \textbf{0.343} & 17 & 53 \\
\textbf{Other 31} (1,278 test pos.) & Transformer & 0.582 & 0.930 & 0.716 &
1,189 & 89 \\
& \textbf{AND ensemble} & 0.776 & 0.810 & \textbf{0.792} & 1,035 & 243 \\
\bottomrule
\end{longtable}

On common categories the conjunction gains +0.077 F1. \textbf{On the rare tail
it nets exactly zero} --- 32.5 points of precision bought with 25.7 points of
recall, detections cut from 35/70 to 17/70. All of the ensemble's aggregate
benefit comes from common categories. \emph{Most Favored Nation} sits in this
tail, is a \textbf{High-risk} category, and scores 0.000 --- which matters more
for a signer than the aggregate does.

\subsection{Step 6 --- Training Stage 2B: Span extractor}

\begin{itemize}
\tightlist
\item
  SQuAD-style \texttt{DistilBertForQuestionAnswering}. Char offsets → token
  start/end via the tokenizer's offset mapping; a 1,200-char focused window is
  re-centered on each answer so it fits in 320 tokens.
\item
  \textbf{11,156 train / 2,667 test} examples; trained on \textbf{8,000}; 2
  epochs, batch 16, LR 3e-5; \textbf{\textasciitilde23 min} MPS.
\end{itemize}

\subsection{Step 7 --- Training Stage 1: Summarizer}

\begin{itemize}
\tightlist
\item
  Clause → one plain-English sentence. CUAD has no plain-language targets, so we
  built a \textbf{41-template seed} (one hand-written gold summary per category,
  weaker-party framing).
\item
  FLAN-T5-small, \textbf{11,156 pairs} (trained on 4,000); 2 epochs, batch 8, LR
  3e-4.
\item
  Training loss 0.141 → 0.085; validation loss 0.096 → 0.075;
  \textbf{\textasciitilde8 min on CPU} (\texttt{use\_cpu=True} to avoid MPS
  OOM).
\end{itemize}

\begin{center}\rule{0.5\linewidth}{0.5pt}\end{center}

\subsection{Step 8 --- TESTING on the held-out 20\% (the metrics)}

All numbers below are on the \textbf{102 unseen test contracts}.

\subsubsection{8.1 Stage 2A --- Presence: the Confusion Matrix (``matrix
test'')}

AND-Ensemble evaluated on all \textbf{4,182 (contract x category) test
predictions}:

\begin{longtable}[]{@{}lll@{}}
\toprule
& Predicted \textbf{Absent} & Predicted \textbf{Present} \\
\midrule
\endhead
\textbf{Actually Absent} & TN = \textbf{2,523} & FP = \textbf{311} \\
\textbf{Actually Present} & FN = \textbf{296} & TP = \textbf{1,052} \\
\bottomrule
\end{longtable}

\textbf{Metrics derived from the confusion matrix:}

\begin{longtable}[]{@{}lll@{}}
\toprule
Metric & Formula & Value \\
\midrule
\endhead
\textbf{Accuracy} & (TP+TN) / all & \textbf{85.49\%} \\
\textbf{Precision} & TP / (TP+FP) & \textbf{77.18\%} \\
\textbf{Recall (Sensitivity)} & TP / (TP+FN) & \textbf{78.04\%} \\
\textbf{Specificity} & TN / (TN+FP) & \textbf{89.03\%} \\
\textbf{F1 Score} & 2·P·R / (P+R) & \textbf{0.7761} \\
\bottomrule
\end{longtable}

\textbf{Per-category F1 (top of the table):}

\begin{longtable}[]{@{}lllll@{}}
\toprule
Category & Test positives & Precision & Recall & F1 \\
\midrule
\endhead
Document Name & 102 & 1.000 & 1.000 & 1.000 \\
Parties & 101 & 0.990 & 1.000 & 0.995 \\
Governing Law & 90 & 0.956 & 0.967 & 0.961 \\
Agreement Date & 93 & 0.929 & 0.989 & 0.958 \\
Expiration Date & 84 & 0.914 & 0.881 & 0.897 \\
License Grant & 56 & 0.907 & 0.875 & 0.891 \\
Anti-Assignment & 74 & 0.831 & 0.932 & 0.879 \\
Effective Date & 73 & 0.831 & 0.877 & 0.853 \\
Cap On Liability & 62 & 0.879 & 0.823 & 0.850 \\
\bottomrule
\end{longtable}

Categories with \textless=7 test positives (e.g.~Source Code Escrow, Most
Favored Nation) score near 0 --- too few examples to learn from in a
510-contract dataset, not a model failure.

\textbf{Why the ensemble is worth keeping despite §5.5.} Precision (77.18\%) and
recall (78.04\%) come out balanced, which is exactly what the AND rule was for:
raw max-pooling trades precision away, and requiring agreement wins it back
without costing recall. Specificity of 89.03\% means the system rarely claims a
clause that is not there --- in a legal-review assistant that is the more
damaging error.

\subsubsection{8.2 Stage 2B --- Span extractor}

Evaluated on \textbf{2,667 answer-bearing test windows}:

\begin{longtable}[]{@{}llll@{}}
\toprule
Metric & Score & 95\% CI (by contract) & 95\% CI (by clause) \\
\midrule
\endhead
\textbf{token-F1} & \textbf{0.763} & {[}0.740, 0.782{]} & {[}0.751, 0.774{]} \\
Overlap (F1 \textgreater= 0.5) & 82.7\% & {[}79.7\%, 85.1\%{]} & {[}81.3\%,
84.1\%{]} \\
Exact Match & 35.5\% & {[}32.5\%, 38.6\%{]} & {[}33.7\%, 37.3\%{]} \\
\bottomrule
\end{longtable}

When the clause is in front of it, the model captures most of the right text
\textasciitilde8 times in 10. The same cluster bootstrap used for presence is
applied here (\texttt{analysis/span\_ci.py}); resampling contracts gives an
interval \textasciitilde2x wider than treating clauses as independent.
\textbf{These numbers assume an answer-bearing window} --- Step 8.4 measures
what happens when Stage 2A picks the window.

\subsubsection{8.3 Stage 1 --- Summarizer}

The 41 targets were \textbf{written by us}, not annotated. Scoring the same
\textbf{150 test clauses} against a second, clause-specific reference set
separates template reproduction from summarization
(\texttt{analysis/summ\_pilot.py}):

\begin{longtable}[]{@{}llll@{}}
\toprule
System & Reference set & ROUGE-L & 95\% CI \\
\midrule
\endhead
Fine-tuned & 41 category templates & \textbf{0.775} & {[}0.718, 0.833{]} \\
Fine-tuned & clause-specific (pilot) & \textbf{0.116} & {[}0.105, 0.128{]} \\
Zero-shot & clause-specific (pilot) & 0.299 & {[}0.275, 0.322{]} \\
\emph{Template alone, no model} & clause-specific (pilot) & 0.115 & {[}0.104,
0.127{]} \\
\bottomrule
\end{longtable}

The fine-tuned model emits one of the 41 templates \textbf{verbatim on 100\% of
clauses} (the correct one 70\% of the time). The bare template with no model
scores the same 0.115 --- the model adds nothing measurable over a lookup. This
is a \textbf{classification-to-template} system and is described as such
throughout; it is not free-text summarization. Pilot references are LLM-drafted
(independent of the model, not gold annotation).

\begin{center}\rule{0.5\linewidth}{0.5pt}\end{center}

\subsubsection{8.4 End-to-end --- what actually reaches the user}

Stage-wise scores do not compose. Running all three stages in series over the
1,348 gold clauses in the 102 test contracts (\texttt{analysis/e2e.py}):

\begin{longtable}[]{@{}llll@{}}
\toprule
Stage & Condition & Surviving & Share \\
\midrule
\endhead
--- & gold clauses & 1,348 & 100.0\% \\
2A Presence & category detected & 1,052 & 78.0\% \\
2B Span & overlap with gold at F1 \textgreater= 0.5 & 293 & 21.7\% \\
Risk & correct level & 293 & 21.7\% \\
\textbf{End-to-end} & & \textbf{293} & \textbf{21.7\%} (CI {[}19.6\%,
23.9\%{]}) \\
\bottomrule
\end{longtable}

Component multiplication predicts \textasciitilde64\%; the pipeline delivers a
third of it. Of the 1,052 detected clauses, the selected window contains the
gold text only \textbf{68.3\%} of the time (max-pooling was trained to detect,
not to locate), and even then extraction succeeds \textbf{39.7\%} vs 82.7\%
isolated --- because training re-centred every focus window on the answer,
narrowing the distribution the model can read. Mean token-F1 falls 0.763 →
0.386. Both causes are fixable without a larger model.

\textbf{For the app:} the risk badge and category are far more reliable than the
quoted clause text.

\begin{center}\rule{0.5\linewidth}{0.5pt}\end{center}

\subsection{Step 9 --- Final scorecard}

\begin{longtable}[]{@{}
  >{\raggedright\arraybackslash}p{(\columnwidth - 10\tabcolsep) * \real{0.17}}
  >{\raggedright\arraybackslash}p{(\columnwidth - 10\tabcolsep) * \real{0.17}}
  >{\raggedright\arraybackslash}p{(\columnwidth - 10\tabcolsep) * \real{0.17}}
  >{\raggedright\arraybackslash}p{(\columnwidth - 10\tabcolsep) * \real{0.17}}
  >{\raggedright\arraybackslash}p{(\columnwidth - 10\tabcolsep) * \real{0.17}}
  >{\raggedright\arraybackslash}p{(\columnwidth - 10\tabcolsep) * \real{0.17}}@{}}
\toprule
Stage & Model & Train & Test & Metric & Score \\
\midrule
\endhead
2A Presence & AND-ensemble (TF-IDF + DistilBERT) & 408 contracts & 102 contracts
& micro-F1 & \textbf{0.776} \\
2A Presence & AND-ensemble & 408 & 102 & Accuracy & \textbf{85.49\%} \\
2A Presence & AND-ensemble & 408 & 102 & Precision / Recall & \textbf{77.18\% /
78.04\%} \\
2B Span & DistilBERT-QA & 408 & 102 & token-F1 (given the window) &
\textbf{0.763} \\
2B Span & DistilBERT-QA & 408 & 102 & token-F1 (window from 2A) &
\textbf{0.386} \\
2B Span & DistilBERT-QA & 408 & 102 & Exact Match & \textbf{35.5\%} \\
1 Summarizer & FLAN-T5-small & 408 & 102 & ROUGE-L (41 templates) &
\textbf{0.775} \\
1 Summarizer & FLAN-T5-small & 408 & 102 & ROUGE-L (clause-specific) &
\textbf{0.116} \\
\textbf{Full pipeline} & all three stages & 408 & 102 & gold clauses surviving &
\textbf{21.7\%} \\
\bottomrule
\end{longtable}

\begin{quote}
\textbf{Read the presence row with §5.5 in hand.} The AND-ensemble's 0.776 is
not statistically separable from the TF-IDF baseline's 0.775 on this test set.
What \emph{is} separable --- and is the result we stand behind --- is that both
beat the windowed transformer alone (0.694) by a wide, significant margin.

\textbf{And read the whole table with §8.4 in hand.} These are component scores
measured with the other stages held out of the way. End to end only
\textbf{21.7\%} of gold clauses survive all three stages. That is the number
describing what a user experiences.
\end{quote}

\begin{center}\rule{0.5\linewidth}{0.5pt}\end{center}

\subsection{Step 10 --- Engineering notes (laptop reality)}

\begin{itemize}
\tightlist
\item
  \textbf{Apple MPS, no CUDA} (\texttt{torch\ 2.8}, \texttt{transformers\ 5.x}).
  DistilBERT + FLAN-T5-small throughout (DeBERTa-v3 is numerically unstable on
  MPS). On a CUDA GPU, swap in DeBERTa-v3-base / flan-t5-base.
\item
  \textbf{MPS memory does not release in-process} → summarizer trained on CPU
  with \texttt{use\_cpu=True}.
\item
  \textbf{\texttt{transformers\ 5.x} API}: \texttt{processing\_class=\ldots{}}
  (not \texttt{tokenizer=\ldots{}}); \texttt{fp16/bf16=False};
  \texttt{dataloader\_pin\_memory=False}.
\item
  \textbf{Leakage discipline} held throughout: contract-level split; retrieval
  queries use only the category/question; gold offsets label training windows
  but never build the input or pick the inference window; the ensemble is
  parameter-free.
\item
  \textbf{One thing we did not keep.} The intermediate checkpoint directories
  were cleared after training, so per-step \textbf{gradient-norm traces were not
  retained} --- only the final weights and the loss values above. The
  learning-rate schedules can be reconstructed exactly (from the saved
  \texttt{training\_args.bin}) and are plotted, but the gradient-norm diagnostic
  would need a re-run (\textasciitilde90 min). Noted as an omission rather than
  glossed over.
\end{itemize}

\subsection{Files}

\begin{verbatim}
clausify-research/
|--- README.md                  — start here: what this is and where to look
|--- thesis/                    — the B.Sc. thesis (LaTeX source + compiled PDF)
|    |--- main.tex / main.pdf    — 55 pages, 8 chapters + appendices
|    |--- chapters/              — one .tex per chapter
|    `--- figures/               — generated plots (TikZ diagrams live in the chapters)
|--- PROJECT_REPORT.md / .pdf   — this complete step-by-step report
|--- REPORT.md / .pdf           — findings-focused report (decisions & caveats)
|--- METHODOLOGY.md / .pdf      — design-decision record
|--- analysis/                  — scripts that generate the reported statistics
|    |--- bootstrap_ci.py        — cluster bootstrap CI (§5.5) + point estimates
|    |--- extra_ci.py            — the three paired model contrasts
|    |--- ch3_chart.py           — train-vs-test split distribution figure
|    |--- lr_schedule.py         — learning-rate schedule figure
|    `--- *.json                 — saved numeric results
|--- data/
|    |--- train_80.csv           — 80% training split (16,728 rows, 408 contracts)
|    `--- test_20.csv            — 20% test split (4,182 rows, 102 contracts)
`--- notebooks/
    |--- README.md              — cell-by-cell notebook guide
    |--- train.ipynb            — training half
    |--- test.ipynb             — testing half (regenerates every number above)
    |--- eda.ipynb              — the original combined notebook
    |--- artifacts/             — the split + the fitted TF-IDF baseline
    `--- outputs/               — the three trained model checkpoints
\end{verbatim}

\textbf{Reproducing the numbers.} \texttt{notebooks/test.ipynb} loads the saved
models from disk and regenerates every metric in this report. The bootstrap
intervals in §5.5 come from \texttt{analysis/bootstrap\_ci.py}, which
independently reproduces the same point estimates (0.776 / 0.775 / 0.718 / 0.696
/ 0.694) and the same confusion matrix (TN=2523, FP=311, FN=296, TP=1052) before
resampling.

\end{document}
